%========================================
% MatinBook - Stage 05: Math System Test
% Test: Equations + Matrices + Operators + Delimiters + Sets
% Purpose: Validate all mathematical features
%========================================

% Add parent directory to search path
\makeatletter
\def\input@path{{../}}
\makeatother

\documentclass{matinbook}

\title{آزمون سیستم ریاضی}
\author{تیم توسعه MatinBook}
\date{تابستان ۱۴۰۵}

\begin{document}
    
    %----------------------------------------
    % Title Page
    %----------------------------------------
    \maketitle
    
    %----------------------------------------
    % Chapter 1: Basic Equations
    %----------------------------------------
    \chapter{معادلات پایه}
    
    \section{معادلات درون خطی}
    
    ریاضیات درون خطی برای اشاره به نمادها و عبارات کوتاه
    استفاده می‌شود: $E = mc^{2}$، قضیه فیثاغورث
    $a^{2} + b^{2} = c^{2}$، و تابع نمایی $e^{i\pi} + 1 = 0$.
    
    همچنین می‌توان از مجموعه‌های عددی استاندارد استفاده کرد:
    $\N$ (طبیعی)، $\Z$ (صحیح)، $\Q$ (گویا)،
    $\R$ (حقیقی)، $\C$ (مختلط).
    
    \section{معادلات شماره‌گذاری شده}
    
    \begin{equation}
        f(x) = x^{3} - 3x^{2} + 2x - 1
        \label{eq:cubic}
    \end{equation}
    
    \begin{equation}
        \frac{d}{dx}\left( \sin x \right) = \cos x
        \label{eq:derivative-sin}
    \end{equation}
    
    معادله \meqref{eq:cubic} یک چندجمله‌ای درجه سه است،
    و معادله \meqref{eq:derivative-sin} مشتق تابع سینوس
    را نشان می‌دهد.
    
    \section{معادلات چندخطی}
    
    محیط \texttt{align} برای معادلات هم‌تراز:
    
    \begin{align}
        (a + b)^{2} &= (a + b)(a + b) \nonumber \\
        &= a^{2} + ab + ba + b^{2} \nonumber \\
        &= a^{2} + 2ab + b^{2} \label{eq:binomial}
    \end{align}
    
    محیط \texttt{gather} برای معادلات غیرهم‌تراز:
    
    \begin{gather}
        x = a + b + c + d + e + f \label{eq:long1} \\
        y = \alpha + \beta + \gamma + \delta + \epsilon \label{eq:long2}
    \end{gather}
    
    محیط \texttt{align} با چند ستون:
    
    \begin{align}
        x &= a + b & y &= c + d \label{eq:multi1} \\
        u &= e + f & v &= g + h \label{eq:multi2}
    \end{align}
    
    %----------------------------------------
    % Chapter 2: Matrices and Vectors
    %----------------------------------------
    \chapter{ماتریس‌ها و بردارها}
    
    \section{انواع ماتریس}
    
    ماتریس با کروشه (\texttt{\textbackslash mat}):
    
    \[
    A = \mat{1 & 2 & 3 \\ 4 & 5 & 6 \\ 7 & 8 & 9}
    \]
    
    ماتریس با پرانتز (\texttt{\textbackslash pmat}):
    
    \[
    B = \pmat{a_{11} & a_{12} \\ a_{21} & a_{22}}
    \]
    
    دترمینان (\texttt{\textbackslash vmat}):
    
    \[
    \det(B) = \vmat{a_{11} & a_{12} \\ a_{21} & a_{22}}
    = a_{11}a_{22} - a_{12}a_{21}
    \]
    
    ماتریس همانی $3 \times 3$:
    
    \[
    I_{3} = \mat{1 & 0 & 0 \\ 0 & 1 & 0 \\ 0 & 0 & 1}
    \]
    
    \section{عملگرهای ماتریسی}
    
    \begin{itemize}
        \item \textbf{اثر (\lr{Trace}):}
        $\trace(A) = a_{11} + a_{22} + \cdots + a_{nn}$
        \item \textbf{دترمینان:} $\det(A) = |A|$
        \item \textbf{ترانهاده:} $A^{T} = [a_{ji}]$
        \item \textbf{ماتریس قطری:}
        $\diag(d_{1}, d_{2}, \ldots, d_{n})$
        \item \textbf{رتبه:} $\rank(A)$
    \end{itemize}
    
    \section{بردارها}
    
    \[
    \vec{v} = \mat{v_{1} \\ v_{2} \\ v_{3}}, \quad
    \vec{w} = \pmat{w_{1} & w_{2} & w_{3}}
    \]
    
    ضرب داخلی: $\inner{\vec{u}}{\vec{v}} = u_{1}v_{1} + u_{2}v_{2} + u_{3}v_{3}$
    
    نرم اقلیدسی: $\norm{\vec{v}} = \sqrt{\inner{\vec{v}}{\vec{v}}}$
    
    %----------------------------------------
    % Chapter 3: Delimiters
    %----------------------------------------
    \chapter{دلیمترهای هوشمند}
    
    \section{قدر مطلق و نرم}
    
    قدر مطلق ساده: $\abs{-5} = 5$
    
    قدر مطلق با اندازه دستی: $\abs[\Big]{\frac{a}{b}}$
    
    نرم بردار: $\norm{\vec{x}}$
    
    نرم با اندیس: $\norm{x}_{2} = \sqrt{\sum_{i} x_{i}^{2}}$
    
    \section{سقف و کف}
    
    \[
    \ceil{3.14} = 4, \quad
    \floor{3.14} = 3, \quad
    \ceil[Big]{\frac{22}{7}} = 4
    \]
    
    \section{ضرب داخلی}
    
    \[
    \inner{u}{v}, \quad
    \inner[\Big]{\frac{x}{2}}{\frac{y}{3}}
    \]
    
    \section{مجموعه}
    
    \[
    A = \set{1, 2, 3, 4, 5}, \quad
    B = \setbuilder{x \in \R}{x > 0}
    \]
    
    \section{پرانتز هوشمند}
    
    \[
    \paren{x + y}, \quad
    \paren[\Big]{\frac{a + b}{c + d}}, \quad
    \paren*{\frac{\sum_{i=1}^{n} x_{i}}{n}}
    \]
    
    %----------------------------------------
    % Chapter 4: Operators and Functions
    %----------------------------------------
    \chapter{عملگرها و توابع ویژه}
    
    \section{عملگرهای حسابان برداری}
    
    گرادیان تابع اسکالر $f$:
    \[
    \grad f = \pderiv{f}{x}\hat{i} + \pderiv{f}{y}\hat{j} + \pderiv{f}{z}\hat{k}
    \]
    
    کرل میدان برداری $\vec{F}$:
    \[
    \curl \vec{F} = \nabla \times \vec{F}
    \]
    
    دیورژانس میدان برداری $\vec{F}$:
    \[
    \diver \vec{F} = \nabla \cdot \vec{F}
    \]
    
    \section{عملگرهای جبر خطی}
    
    \begin{align}
        \rank(A) &= \text{بعد فضای ستونی} \label{eq:rank} \\
        \trace(A) &= \sum_{i=1}^{n} a_{ii} \label{eq:trace} \\
        \diag(d_{1}, \ldots, d_{n}) &=
        \mat{d_{1} & 0 & \cdots & 0 \\
            0 & d_{2} & \cdots & 0 \\
            \vdots & \vdots & \ddots & \vdots \\
            0 & 0 & \cdots & d_{n}} \label{eq:diag}
    \end{align}
    
    \section{عملگرهای نظریه اعداد}
    
    \begin{itemize}
        \item بزرگترین مقسوم‌علیه مشترک: $\gcdop(a, b)$
        \item کوچکترین مضرب مشترک: $\lcm(a, b)$
        \item تابع علامت: $\sgn(x)$
    \end{itemize}
    
    \section{بهینه‌سازی}
    
    \begin{align}
        \argmax_{x \in \R} f(x) &= \text{مقدار } x \text{ که } f(x) \text{ را بیشینه کند} \\
        \argmin_{x \in \R} f(x) &= \text{مقدار } x \text{ که } f(x) \text{ را کمینه کند}
    \end{align}
    
    %----------------------------------------
    % Chapter 5: Calculus
    %----------------------------------------
    \chapter{حساب دیفرانسیل و انتگرال}
    
    \section{مشتق}
    
    مشتق معمولی:
    \[
    \deriv{}{x}\left( x^{n} \right) = n x^{n-1}
    \]
    
    مشتق جزئی:
    \[
    \pderiv{f}{x} = \lim_{h \to 0} \frac{f(x+h, y) - f(x, y)}{h}
    \]
    
    مشتق دوم:
    \[
    \deriv{^{2}f}{x^{2}} = \deriv{}{x}\left( \deriv{f}{x} \right)
    \]
    
    \section{انتگرال}
    
    انتگرال نامعین:
    \[
    \int x^{n} \dd{x} = \frac{x^{n+1}}{n+1} + C
    \]
    
    انتگرال معین:
    \[
    \int_{0}^{\infty} e^{-x^{2}} \dd{x} = \frac{\sqrt{\pi}}{2}
    \]
    
    انتگرال دوگانه:
    \[
    \iint_{D} f(x, y) \dd{x}\dd{y}
    \]
    
    \section{حد و سری}
    
    حد معروف:
    \[
    \lim_{n \to \infty} \left(1 + \frac{1}{n}\right)^{n} = e
    \]
    
    سری تیلور:
    \[
    e^{x} = \sum_{n=0}^{\infty} \frac{x^{n}}{n!}
    = 1 + x + \frac{x^{2}}{2!} + \frac{x^{3}}{3!} + \cdots
    \]
    
    سری فوریه:
    \[
    f(x) = \frac{a_{0}}{2} + \sum_{n=1}^{\infty}
    \left[ a_{n} \cos(nx) + b_{n} \sin(nx) \right]
    \]
    
    %----------------------------------------
    % Chapter 6: Special Functions
    %----------------------------------------
    \chapter{توابع ویژه}
    
    \section{تابع گاما}
    
    \[
    \Gamma(z) = \int_{0}^{\infty} t^{z-1} e^{-t} \dd{t}, \quad \Re(z) > 0
    \]
    
    \section{تابع سینک}
    
    \[
    \sinc(x) = \frac{\sin(\pi x)}{\pi x}
    \]
    
    \section{تابع خطا}
    
    \[
    \erf(x) = \frac{2}{\sqrt{\pi}} \int_{0}^{x} e^{-t^{2}} \dd{t}
    \]
    
    %----------------------------------------
    % Chapter 7: Number Sets
    %----------------------------------------
    \chapter{مجموعه‌های عددی}
    
    \section{نمادهای استاندارد}
    
    \begin{table}[h]
        \centering
        \caption{مجموعه‌های عددی استاندارد}
        \label{tab:number-sets}
        \begin{tabular}{|c|c|C{7cm}|}
            \hline
            \textbf{نماد} & \textbf{نام} & \textbf{توضیحات} \\
            \hline
            $\N$ & اعداد طبیعی & $\{1, 2, 3, \ldots\}$ \\
            $\Z$ & اعداد صحیح & $\{\ldots, -2, -1, 0, 1, 2, \ldots\}$ \\
            $\Q$ & اعداد گویا & $\{p/q \mid p,q \in \Z, q \neq 0\}$ \\
            $\C$ & اعداد مختلط & $\{a + bi \mid a,b \in \R\}$ \\    
            \hline
        \end{tabular}
    \end{table}
    
    \section{رابطه بین مجموعه‌ها}
    
    \[
    \N \subset \Z \subset \Q \subset \R \subset \C
    \]
    
    %----------------------------------------
    % Summary
    %----------------------------------------
    \chapter{خلاصه نتایج}
    
    \section{وضعیت ماژول ریاضی}
    
    \begin{table}[h]
        \centering
        \caption{وضعیت قابلیت‌های ریاضی}
        \label{tab:math-status}
        \begin{tabular}{|c|C{6cm}|c|}
            \hline
            \textbf{بخش} & \textbf{قابلیت‌ها} & \textbf{وضعیت} \\
            \hline
            معادلات & \lr{align, gather, equation} & \textcolor{matingreen}{✅} \\
            ماتریس‌ها & \lr{mat, pmat, vmat} & \textcolor{matingreen}{✅} \\
            دلیمترها & \lr{abs, norm, ceil, floor, inner, set, paren} & \textcolor{matingreen}{✅} \\
            عملگرها & \lr{grad, curl, diver, rank, trace, diag} & \textcolor{matingreen}{✅} \\
            مشتق/انتگرال & \lr{deriv, pderiv, dd, D} & \textcolor{matingreen}{✅} \\
            مجموعه‌ها & \lr{N, Z, Q, R, C} & \textcolor{matingreen}{✅} \\
            توابع ویژه & \lr{sinc, sgn, gcdop, lcm, erf} & \textcolor{matingreen}{✅} \\
            \hline
        \end{tabular}
    \end{table}
    
    \section{نتیجه}
    
    \textbf{ماژول ریاضی \lr{MatinBook} نسخه ۱ با موفقیت آزمایش شد.}
    تمام معادلات، ماتریس‌ها، دلیمترها، عملگرها و مجموعه‌های عددی
    به درستی کار می‌کنند. 🎉
    
\end{document}